\documentclass[11pt,a4paper]{article}
\usepackage[utf8]{inputenc}
\usepackage[T1]{fontenc}
\usepackage{lmodern}
\usepackage[french]{babel}
\usepackage{amsmath,amssymb,mathtools}
\usepackage{icomma}
\usepackage{xcolor}
\usepackage{tikz}
\usepackage{pgfplots}
\usepackage[most]{tcolorbox}
\usepackage{enumitem}
\usepackage{fancyhdr}
\usepackage[hidelinks]{hyperref}
\usepackage{titlesec}
\usepackage[a4paper,margin=2.2cm,headheight=26pt,headsep=10pt]{geometry}
\usetikzlibrary{arrows.meta,positioning,calc,intersections,angles,quotes}
\usepgfplotslibrary{fillbetween}
\pgfplotsset{compat=1.18}
\definecolor{primary}{RGB}{30,75,145}
\definecolor{primarydark}{RGB}{20,50,100}
\definecolor{defcol}{RGB}{0,114,178}
\definecolor{thmcol}{RGB}{0,135,110}
\definecolor{excol}{RGB}{0,130,85}
\definecolor{remarkcol}{RGB}{120,70,180}
\definecolor{attncol}{RGB}{200,40,55}
\definecolor{methcol}{RGB}{85,85,100}
\definecolor{areacol}{RGB}{120,160,230}
\definecolor{areacol2}{RGB}{230,150,80}
\newcommand{\dd}{\mathrm{d}}
\newcommand{\R}{\mathbb{R}}
\newcommand{\ua}{\,\text{u.a.}}
\newcommand{\tg}{\tan}\newcommand{\cotg}{\cot}
\tcbset{boxbase/.style={enhanced,breakable,boxrule=0.9pt,arc=2.5pt,left=3mm,right=3mm,top=2.3mm,bottom=2.3mm,fonttitle=\bfseries,coltitle=white,toptitle=0.6mm,bottomtitle=0.6mm}}
\newtcolorbox{methode}[1][]{boxbase,colback=methcol!7,colframe=methcol,sharp corners,title={\textsc{Comment faire}\ifx\relax#1\relax\relax\else\textmd{ : \itshape #1}\fi}}
\newtcolorbox{exemple}[1][]{boxbase,colback=excol!5,colframe=excol,title={\textsc{Exemple}\ifx\relax#1\relax\relax\else\textmd{ : \itshape #1}\fi}}
\newtcolorbox{idee}[1][]{boxbase,colback=defcol!5,colframe=defcol,title={\textsc{Idée clé}\ifx\relax#1\relax\relax\else\textmd{ : \itshape #1}\fi}}
\newtcolorbox{remarque}[1][]{boxbase,colback=remarkcol!6,colframe=remarkcol,title={\textsc{Remarque}\ifx\relax#1\relax\relax\else\textmd{ : \itshape #1}\fi}}
\newtcolorbox{attention}[1][]{boxbase,colback=attncol!6,colframe=attncol,title={\textsc{Attention}\ifx\relax#1\relax\relax\else\textmd{ : \itshape #1}\fi}}
\newtcolorbox{exo}[1][]{boxbase,colback={primary!4},colframe=primary,title={\textsc{Exercice #1}}}
\newtcolorbox{corr}[1][]{boxbase,colback={thmcol!5},colframe=thmcol,title={\textsc{Corrigé #1}}}
\tikzset{axe/.style={->,>=Stealth,black!65,line width=0.7pt},courbe/.style={primary,line width=1.3pt},courbe2/.style={areacol2!90!black,line width=1.3pt},dvert/.style={gray,dashed,line width=0.7pt}}
\pagestyle{fancy}\fancyhf{}
\fancyhead[L]{\small\textcolor{primarydark}{\textbf{Fiche 7} — Intégrale d'une fonction continue}}
\fancyhead[R]{\small\textcolor{primarydark}{\textsc{Thème 5 — Calculs d'aires}}}
\fancyfoot[C]{\small\thepage}
\renewcommand{\headrulewidth}{0.8pt}
\titleformat{\section}{\large\bfseries\color{primarydark}}{\thesection}{0.6em}{}

\begin{document}

\section*{Exercices — Difficile : aire entre deux courbes}

\begin{center}\itshape\small
Méthode : (1) bornes $=$ solutions de $f(x)=g(x)$ ; (2) déterminer la courbe supérieure ;
(3) intégrer haut $-$ bas ; (4) calculer et vérifier la positivité.
\end{center}

\begin{exo}[D1 — deux paraboles]
On considère $f(x)=x^2-4x+3$ et $g(x)=-x^2+4x-3$.
\begin{enumerate}[label=\alph*),leftmargin=2em,topsep=2pt,itemsep=1pt]
  \item Déterminer les abscisses des points d'intersection de $\mathcal{C}_f$ et $\mathcal{C}_g$.
  \item Préciser, sur cet intervalle, laquelle des deux courbes est au-dessus.
  \item Écrire l'intégrale donnant l'aire $\mathcal{A}$ du domaine compris entre les deux courbes.
  \item Calculer $\mathcal{A}$.
\end{enumerate}

\smallskip
\begin{center}
\begin{tikzpicture}
  \begin{axis}[width=8.8cm,height=4.6cm,axis lines=middle,
      xmin=-0.4,xmax=4.4,ymin=-1.8,ymax=1.8,xtick={1,3},ytick=\empty,
      xlabel={$x$},ylabel={$y$},
      every axis x label/.style={at={(current axis.right of origin)},anchor=west},
      every axis y label/.style={at={(current axis.above origin)},anchor=south}]
    \addplot[domain=0.2:3.8,samples=80,courbe]{x^2-4*x+3};
    \addplot[domain=0.2:3.8,samples=80,courbe2]{-x^2+4*x-3};
    \addplot[name path=G,domain=1:3,samples=50,draw=none,forget plot]{-x^2+4*x-3};
    \addplot[name path=F,domain=1:3,samples=50,draw=none,forget plot]{x^2-4*x+3};
    \addplot[areacol!40] fill between[of=G and F];
    \node[primary] at (axis cs:3.7,1.1){$f$};
    \node[areacol2!90!black] at (axis cs:3.7,-1.15){$g$};
    \node at (axis cs:2,0){$\mathcal{A}$};
  \end{axis}
\end{tikzpicture}
\end{center}
\end{exo}

\begin{exo}[D2 — parabole et droite]
On considère la parabole $f(x)=x^2-4x+3$ et la droite $g(x)=x-1$.
\begin{enumerate}[label=\alph*),leftmargin=2em,topsep=2pt,itemsep=1pt]
  \item Déterminer les abscisses des points d'intersection de $\mathcal{C}_f$ et de la droite.
  \item Déterminer l'ordre des deux courbes sur l'intervalle obtenu.
  \item Calculer l'aire $\mathcal{A}$ du domaine délimité par la parabole et la droite.
\end{enumerate}
\end{exo}

\begin{exo}[D3 — deux paraboles]
On considère $f(x)=x^2-4x$ et $g(x)=2x-x^2$.
\begin{enumerate}[label=\alph*),leftmargin=2em,topsep=2pt,itemsep=1pt]
  \item Résoudre $f(x)=g(x)$ pour trouver les bornes.
  \item Déterminer la courbe supérieure sur cet intervalle.
  \item Calculer l'aire $\mathcal{A}$ comprise entre les deux paraboles.
\end{enumerate}

\smallskip
\begin{center}
\begin{tikzpicture}
  \begin{axis}[width=8.8cm,height=4.6cm,axis lines=middle,
      xmin=-1,xmax=4.2,ymin=-4.6,ymax=1.6,xtick={0,3},ytick=\empty,
      xlabel={$x$},ylabel={$y$},
      every axis x label/.style={at={(current axis.right of origin)},anchor=west},
      every axis y label/.style={at={(current axis.above origin)},anchor=south}]
    \addplot[domain=-0.7:4,samples=80,courbe]{x^2-4*x};
    \addplot[domain=-0.7:4,samples=80,courbe2]{2*x-x^2};
    \addplot[name path=G,domain=0:3,samples=50,draw=none,forget plot]{2*x-x^2};
    \addplot[name path=F,domain=0:3,samples=50,draw=none,forget plot]{x^2-4*x};
    \addplot[areacol!40] fill between[of=G and F];
    \node[primary] at (axis cs:3.9,-0.3){$f$};
    \node[areacol2!90!black] at (axis cs:1,1.15){$g$};
    \node at (axis cs:1.5,-1.3){$\mathcal{A}$};
  \end{axis}
\end{tikzpicture}
\end{center}
\end{exo}

\begin{exo}[D4 — courbes qui se croisent dans l'intervalle]
On considère $f(x)=x^2$ et $g(x)=x$ sur l'intervalle $[0,2]$.
\begin{enumerate}[label=\alph*),leftmargin=2em,topsep=2pt,itemsep=1pt]
  \item Résoudre $f(x)=g(x)$. Montrer que les deux courbes se croisent \emph{à l'intérieur} de $[0,2]$.
  \item Déterminer, sur $[0,1]$ puis sur $[1,2]$, laquelle des deux courbes est au-dessus.
  \item En déduire l'aire totale $\mathcal{A}$ comprise entre $\mathcal{C}_f$ et $\mathcal{C}_g$ sur $[0,2]$
        (on coupera en $x=1$).
\end{enumerate}
\end{exo}

\end{document}
